{
\setlength{\parindent}{2em}
\chapter{Flight Software Checkpointing}\label{cha:das-impl}
At its core, offering a suitable environment for the simulation of the flight software was the main purpose of the BBPSim simulator. Since \gls{SNC} wanted to reproduce the in-flight behavior of the communication sub-system as faithfully as possible, the software-in-the-loop test framework was designed according to the principle of encapsulation: exclude all internal changes to the simulated software and instead build around what is already there. Having said that, inevitably, there were exceptions made for some parts of the flight software algorithms, where code had to be altered to make it compatible to run on a Linux machine. For the most part, however, as shown in \autoref{fig:bbpsim-layers}, BBPSim succeeded in integrating \gls{DAS} as an "untouchable" nucleus and interfacing only indirectly with it.

The fact that DAS was supposed to stay completely permeable to modification forced the implementation of the save \& restore feature to treat it as a black box. This point of view brought several challenges, which were not present when modifying the BBPSim code to support checkpointing. 

In this section, the constraints that influenced the development of a checkpointing technique within the flight software are first of all described. It should be noted that most of the solutions to the problems were based on technical details that, in the end, had a meaningful impact. Subsequently, \gls{DAS} is briefly analyzed to clearly define which of its components' checkpointing are considered a necessary condition for a stable restore. Then, the approach taken to indirectly access, save and restore them is explained in details.

\section{Design Constraints}
It goes without saying that user and customer requirements play a big role in how the design goes forward. In particular, to make the \gls{FSW} checkpointable, it was important to bear in mind DCMS-BBP-121 (64-bit application), U02 (no alteration to FSW) and U01 (no instability). However, some other technicalities had a big importance too.

First and foremost, since the delivery method of BBPSim was through a Linux dynamic library, it was realized that most of the container-based solutions presented in \autoref{cha:state-of-the-art} could not apply to BBPSim. A shared object does not possess the same control than an application over many things, like memory layout. This presented a big problem, and a better explanation is given in the coming sections.

TO ADD HERE 

\section{Definition of Sufficient Condition for Restoration}\label{sec:conditions}
Saving and restoring an application was the main subject of \autoref{cha:state-of-the-art}. In every existing solution that was presented, the checkpointer programs did not have any control over the building process of the checkpointee. This made them approach the checkpointing problem from various angles, depending on the level at which the program was supposed to interact with the checkpointed application. 

In parallel, they each introduced a different sufficient condition, in terms of what to save, to be able to restore back the checkpointee. For example, it was considered essential for CRIU to save the totality of active memory pages inside the application's \gls{VMA} range, something pointless using the C\textsuperscript{3} environment.

Since the checkpointing feature of BBPSim was done with a different set of constraints, the set of necessary conditions for a stable restore had to once again be revisited. In this section, an assessment of the required components to checkpoint in the flight software is done. 

\subsection*{Definition of FSW Program State}
The description of requirement DCMS-BBP-122 on page \pageref{tab:customer-reqs} mentioned that "the current state of the [BBPSim] model" had to be saveable between steps. Undoubtedly, the word "model" is a very abstract term, and it could be interpreted in many ways. It was important to first define what qualified as a "model" in the more concrete context of DAS.

To better picture this, a parallel can be drawn with how C code is structured once compiled. When the \gls{GCC} compiles a valid C source file on Linux, an object file with extension \texttt{.o} is produced. Internally, this file is laid out in the Unix-wide standard \gls{ELF}. This very flexible format is divided into multiple segments that each have their own function in defining the program at low-level.

In this context, there are four relevant memory segments inside an object file:
\begin{wrapfigure}{r}{0.45\textwidth}
	\centering 
	\includegraphics[width=.9\linewidth,keepaspectratio]{art/reloc-obj.png}
	\caption{Sections in an object file with ELF format.\cite{online:zhang}}
	\label{fig:sections-obj}
	\vspace{-24pt}
\end{wrapfigure}
\begin{enumerate}
	\item \textbf{\texttt{.bss}}. Contains the statically-allocated variables that were either uninitialized or initialized to zero.
	\item \textbf{\texttt{.data}}. Contains the statically-allocated variables that were initialized to zero.
	\item \textbf{\texttt{.text}}. Contains the binary instructions. Basically, the executable code.
	\item \textbf{\texttt{.rodata}}. Contains the read-only data (constants).
\end{enumerate}
\autoref{code:c-to-segments} shows how the translation is done using real C code.

In this thesis, it is argued that the state, or "model" of the flight software could be defined by two main components: the memory segments of statically-allocated variables, and the program counters. In the context of an object file, it was possible to redefine the state of a multithreaded DAS as being:
\begin{shadedquotation}
The full contents of the \texttt{.bss} and \texttt{.data} segments for all object files, the register set of each threads and the stack of each threads. 
\end{shadedquotation}

As for the \texttt{.rodata} and other segments, since they contain values that do not change from one simulation to the other, they were deemed unnecessary to include in this definition. Since they were identical at every simulation run, they were not included in the definition of the FSW state.

\subsection*{Condition 1 - Memory Segments}
The first necessary condition to be able to restore back to a previous "model" of DAS was to save the read-writeable, \textit{statically-allocated} regions in memory that concerned the flight software. It should be understood that a program's statically-allocated variables are variables that exist for the entire duration of the program. This term should not be confused with the \mintinline{c}|static| keyword of the C language, which limits the scope of a symbol in a compilation unit.

One can make an informal proof by contradiction of that statement by looking again at the sample code in \autoref{code:c-to-segments}. Let's suppose that, to restore DAS to a consistent global state, the restoration of the statically-allocated memory segments is not necessary. As a user, any changes done to variables (declared on lines \ref{lc-beg-statics} to \ref{lc-beg-statics}) during a simulation will be lost, and therefore, the global state of the application would not be consistent with the one on the previous run. Variables on the second run would not read the same as they would have at the end of the first one. The inclusion of these segments is thus a necessary condition.

\subsection*{Condition 2 - Register Sets}
A second necessary condition for a consistent DAS global state restore was to save the register set of every thread attached to the DAS program. In this set is located the \gls{PC}, depicted in \autoref{fig:x86-regs}. It is the CPU register in charge of keeping track of which instructions are to be executed. Since CPU instructions are encoded in binary and reside in the same \gls{VMA} as the variables, they are also addressable. In that sense, the program counter register holds the address of the current instruction being executed by the processor. One can think of the PC as being "where" the thread is in its execution of the program.

In addition, it was also important to save the other general-purpose registers, which contained the context of the CPU at that particular moment. Not putting the correct value in the registers would produce a domino effect resulting in a segmentation fault at best, or undefined behavior at worst.

\subsection*{Condition 3 - Thread Execution Stacks}
Finally, saving the thread execution stacks was considered the last necessary condition to make the checkpointing feature work in the context of the flight software. An execution stack contains a lot of very important information about the thread. For instance, the function-local (scoped) variables are located in this stack. For reasons similar to the first condition, it was also important to save them.

On top of that, even though the program counter represents the position of the execution of a thread at a certain time $t$, the stack also contains its \textit{backtrace}, information about the "path" taken by the thread to get there. Precisely, the backtrace (or stack trace) is defined as the series of functions that were called consecutively until $t$. The stack-related concepts are treated thoroughly in \autoref{sec:das-exec-restore}.

\subsection*{Building the Sufficient Condition}
When the three conditions listed above were grouped together, it was defined that this formed a sufficient condition for a consistent DAS global state restore. More formally, this thesis argues that:
\begin{itemize}
	\item with $F$ being a C function in DAS;
	\item with $G_s$ being the global state of DAS in memory;
	\item with $S_r$ being a previously saved state being restored;
\end{itemize}
then
\begin{equation} \label{eq:sr_equiv}
	\forall F: F(G_s)\equiv F(S_r)
\end{equation}
when $F$ is void of non-deterministic processes. All DAS functions satisfied to this criteria, because they were strictly exempt of dynamic memory allocation, random numbers, and scheduling non-determinism (since the order of execution for threads is set at compile-time).

It should be noted that \autoref{eq:sr_equiv} also holds true whether $F$ is \textit{stateful} or stateless (i.e whether calling $F$ with the same parameters yields the same output or not). This is because even function-scoped static variables are included in the memory segments that are part of $S_r$. \autoref{code:c-to-segments} contains a good example of this. The \texttt{count()} function is stateful, but its state (i.e the value of \mintinline{c}|theCount|) is also included in the definition of the flight software state (in the \texttt{.bss} section).

Now that \textit{what} to save is properly defined, the following sections detail how the conditions were met for each of them to achieve a consistent checkpointing of the flight software.

%big sections
{
\setlength{\parindent}{2em}
\section{Black Box Symbol Access}\label{sec:das-mem-restore}
It was mentioned previously that, for integration reasons, the flight software code was considered unalterable (through requirement U02). Modifying the \gls{FSW} to accommodate the new checkpointing feature in BBPSim clearly defeated the purpose of encapsulation, and could potentially introduce errors in a critical piece of software. To represent reality as much as possible, the FSW was supposed to contain as fewer indications of the existence of \gls{BBPSim} as possible in its code.

In \autoref{sec:conditions}, it was shown that saving the entire set of statically-allocated variables in the DAS domain was a necessary condition for a stable restore. However, as it was forbidden to add \textit{getter} functions (e.g. \mintinline{c}|GetTheVariable()|) to access file-local variables from BBPSim, there needed to be a different solution. For the save and restore to be properly integrated, there needed to be a way to access them at runtime, a kind of indirect or "black box" access. 

In the end, the solution that was adopted involved the handling of DAS object files, \textit{after} they were compiled. Ultimately, no control over the flight software source code was given, but the compilation and building process were alterable. The following sections show how this limited control could be harvested, using a series of manipulations at build time, to give the \gls{BBPSim}-domain objects the required access to save the state of the flight software.

\subsection*{Linking Process}
First of all, it is crucial to understand how the linking process of an application is done. As mentioned earlier, when each source code file is compiled by \gls{GCC}, a \texttt{.o} object file is produced. Although this file contains executable code, the object file that would be produced by GCC would, by itself, not be enough to be executable as a stand-alone. The \textit{linker} needs to first process the object files. The one used in the case of this thesis was GNU \texttt{ld} version 2.27.

\begin{figure}[htbp]
	\centering 
	\includegraphics[width=.75\linewidth,keepaspectratio]{art/obj-to-elf-to-mem.png}
	\caption{Combining Input Sections to Form an Executable Object Module.\cite{online:linking}}
	\label{fig:link-mult-files}
\end{figure}

The linker's role in the building process is to combine each of the object files given as input into an executable file, or in the case of this thesis, a dynamic library. \autoref{fig:link-mult-files} better illustrates that concept. At linking time, \texttt{ld} groups the same segments from each \texttt{.o} file together and lays them sequentially in the executable object. Since object files can be given in an arbitrary order, the linker can make no guarantee on the final layout on its own.

In addition, while combining the object files together, the linker also resolves the undefined \textit{external} symbols, variables declared in one object file but used in another. 

Just like source code determines the resulting output of the compiler, it is also possible to control the linking process through the use of a linker script\cite{online:ld-scripts}. Even though, in most cases, users of the GNU linker don't provide a custom script, it is possible to specify one at linking time. 

Using such a script offers much more granularity to the developer in terms of which section gets put where in the executable object. In the field of embedded systems, this is particularly useful. Since components like non-volatile flash memory can be adressable by the CPU\cite{online:flash-ram}, it lets the developer decide whether certain sections should be written to flash (e.g. addressed from 0x10000 to 0x50000) or just be allocated in RAM before running the program.

\begin{listing}[H]
	\vspace{12pt}
	\begin{minted}{c}
SECTIONS
{
	. = 0x10000;
	.text : { *(.text) }//$\label{line:link-ex}$
	. = 0x8000000;
	.data : { *(.data) }
	.bss : { *(.bss) }
}
	\end{minted}
	\caption{Simple example of a GNU \texttt{ld} linker script.}
	\label{code:link-script-ex}
\end{listing}

One can read the example linker script in \autoref{code:link-script-ex}  sequentially. On line \ref{line:link-ex}, the script lays out the \texttt{.text} sections of all the object files sequentially inside the executable object's own \texttt{.text} section. This newly created section starts at address 0x10000 (represented by the "." keyword, the \textit{location counter}\cite{online:gnu-ld}). Then, the linker does the same thing for the \texttt{.data} and \texttt{.bss} sections, starting at address 0x8000000.

\subsection*{Clustering Flight Software Sections}
Since managing memory is always easier when it is contiguous, the first step that was taken to access the flight software memory was to group the sections of the relevant files together. This was done using a custom linker script. To make sure not to leave out any section, the default GNU \texttt{ld} linker script was used to make sure to start from a valid point. This script could be outputted from a shell with 
\begin{minted}{bash}
ld --verbose
\end{minted}
From that point, it was possible to partition the essential sections of the DAS object files only by changing the already valid script. These sections could be clustered in a special section of an intermediate shared object \pathmono{libBbpSim-intermediate.so}. 

As previously stated, one of the linker's job when building an executable is to resolve undefined symbols across all object files and to "interlink" them together. To do that, \texttt{ld} starts by looking in other object files to locate and set the correct addresses. Despite that, if said symbols are not defined in any object files, \textit{they can also be defined in the linker script}. This has two major implications: 
\begin{enumerate}
	\item Bi-directional transmission of values is possible (C code $\Longleftrightarrow$ linker).
	\item The source code can never know in advance the value of the undefined symbol if it's only defined during the linking phase.
\end{enumerate}

By harvesting these concepts together, two new contiguous sections of DAS variables in \pathmono{libBbpSim-intermediate.so} could be constructed. They are shown in \autoref{code:link-script-das}. For both the new \texttt{.das_state_vars_bss} and \texttt{.das_state_vars_data} sections, two symbols were defined as boundary variables to identify the start and end of the section in memory. A pair can be seen at lines \ref{line:bound-start} and \ref{line:bound-end}.

\begin{listing}[htbp]
	\vspace{12pt}
	\begin{minted}{c}
//...
.das_state_vars_data :
{
	__das_state_vars_data_start = .;//$\label{line:bound-start}$
	SORT(*obj/DasSrc/ *.o)(.data .data.* .gnu.linkonce.d.*)
	__das_state_vars_data_end = .;//$\label{line:bound-end}$
}
.das_state_vars_bss (NOLOAD) :
{
	__das_state_vars_bss_start = .;
	SORT(*obj/DasSrc/ *.o)(.bss .bss.* COMMON)
	__das_state_vars_bss_end = .;
}
//...
	\end{minted}
	\caption{Clustering of DAS variables via the \texttt{ld} linker script.}
	\label{code:link-script-das}
\end{listing}

It should also be noted that the \texttt{NOLOAD} parameter of the new \texttt{.das_state_vars_bss} section informed the linker that this memory section was to be zeroed out when the library was loaded, and thus not to allocate this memory inside the shared object itself. Otherwise, the library would be composed of nearly 100MB of zeros. And thanks to the inclusion of section guards, it was now possible to access the start and end addresses of the sections from C using:
\begin{minted}{c}
extern char __das_state_vars_data_start[];
extern char __das_state_vars_data_end[];
\end{minted}

To confirm that the script worked, an output map was produced by \texttt{ld} after the linking process (using \texttt{-Wl,-M=mapfile.map}). An excerpt is given in \autoref{map:ld}. It shows that the sections get created and that the \texttt{.data} and \texttt{.bss} sections of each DAS object file all get correctly mapped in the output file.

\subsection*{Cataloging State Variables}
Once the relevant sections of memory related to DAS were correctly laid out after the linking process, it was now a matter of deducing the address of each of these variables inside the resulting shared object. As the outputted file was in the ELF format, an ELF analyzing program could then be used to get the list of all the variables contained within the relevant section:
\begin{minted}{bash}
objdump --syms -j .das_state_vars_data libBbpSim-intermediate.so
\end{minted}
Once the output was sanitized with \texttt{awk}, a big text file represented in  \autoref{das-symbol-catalog} was produced. This was called the \textit{symbol catalog}. One can observe various components in this file:
\begin{enumerate}
	\item The first line contains the SHA1 sum of the entire DAS source code repository. This was added to associate the list of variables to an existing source code for compatibility purposes.
	\item Each additional line represents a symbol that is contained in the flight software domain. They are laid out in the \verb|[<address> <size> <name>]| format. The addresses and the sizes are outputted in hexadecimal.
	\item The symbols defined as "guards" (\texttt{__das_state_vars_data_start}) earlier in the linker script were also included in this list. Because they don't represent any value in source code but instead hold an address, their size was zero.
\end{enumerate}

\subsection*{Including the Catalog in the Library}
Being able to generate a symbol catalog was very useful. However, there was a problem: with the way it was constructed, that catalog could never be included in the library. The catalog was generated from a finished shared object, but for it to ever be useful, the catalog needed to be accessible by the library at runtime. This created a chicken and egg problem, where both artifacts were dependent on one another.

The dependence was bypassed by transforming the symbol catalog into an object file and re-linking the library once again, including the new transformed catalog. 

Embedding raw bytes inside an executable file has been done since the early days of computing. For instance, it is possible for a malicious programmer to embed executable code within a self-extracting archive file\cite{online:sfx}. In the case of BBPSim, which was an ELF running on an x86 64-bit machine, it was possible to embed the catalog, as a string, \textit{inside} the library itself by first transforming it into a \texttt{.o} file:
\begin{minted}{bash}
objcopy --input binary --output elf64-x86-64          \
--binary-architecture i386:x86-64                     \
--rename-section .data=.das_state_vars_catalog_string \
das-symbol-catalog.txt                                
\end{minted}

This put the raw bytes of the catalog in the \texttt{.das_state_vars_catalog_string} section of the \pathmono{das-symbol-catalog.o} file.

The second part consisted in correctly laying out those bytes in the second and final shared object, \pathmono{libBbpSim.so}. This was done by once again taking advantage of the linker script in \autoref{code:link-script-das} to add an additional section, \texttt{.das_state_vars_catalog_string}.

Finally, all the object files both from the BBPSim and DAS domains could be linked a second time, together with \pathmono{das-symbol-catalog.o}. The resulting \pathmono{libBbpSim.so} file would be the last and final shared object. This effectively embedded the symbol catalog's raw bytes, encoded in \gls{ASCII} characters. Executable code inside the library could also refer to the stringized symbol catalog using newly created symbols used as guards for the raw bytes. Those symbols being only defined in the second linking process, it was necessary to tag them as having \texttt{weak} linkage for the linker not to complain on the first linker pass:
\begin{minted}{c}
extern char __attribute__((weak)) _binary_das_symbol_catalog_txt_start[];
extern char __attribute__((weak)) _binary_das_symbol_catalog_txt_end[];
\end{minted}

The results can be better shown by using a binary visualization program\cite{online:binvisio}. \autoref{fig:bin-lib-comp} shows a comparison of the \pathmono{libBbpSim-intermediate.so} and \pathmono{libBbpSim.so} shared objects in binary. Bytes with a value in the \gls{ASCII} character range are represented in blue. Both images read left-to-right with address 0x0 on the left. It is possible to see that a new blue section appears in \autoref{fig:bin-lib-final}, thus confirming that the symbol catalog correctly gets included in the final version of the BBPSim library.

\begin{figure}[htbp]
	\centering
	\begin{subfigure}{\linewidth}
		\centering
		\includegraphics[width=\linewidth, keepaspectratio]{art/bin-lib-intermediate.png}
		\caption{Binary visualization of \pathmono{libBbpSim-intermediate.so}}
		\label{fig:bin-lib-intermediate}
	\end{subfigure}
	\vspace{12pt}
	\begin{subfigure}{\linewidth}
		\centering
		\includegraphics[width=\linewidth]{art/bin-lib-final.png}
		\caption{Binary visualization of \pathmono{libBbpSim.so}}
		\label{fig:bin-lib-final}
	\end{subfigure}
	\caption{Comparison of the resulting libraries after the first and second linking process.}
	\label{fig:bin-lib-comp}
\end{figure}

Since there is no modification to the DAS object files between the two linking process, the variable list taken from \pathmono{libBbpSim-intermediate.so} \textit{is guaranteed} to be the same and in the same order as \pathmono{libBbpSim.so}. 

\subsection*{Shared Object Considerations}
- Dynamic library loading  https://eli.thegreenplace.net/2011/08/25/load-time-relocation-of-shared-libraries
- address at which the library is loaded (why alwayshttps://unix.stackexchange.com/questions/509607/how-a-64-bit-process-virtual-address-space-is-divided-in-linux 0x00007fffXXXXXXXX?) 
In this case, application continerization is not possible, because we are a library.
- how to access everything (talk about dynamic library loading, elf dumps, double linking etc.)

\subsection*{Summary}
The whole process from is quite complex. With this manipulation of sections, one can get easily confused. This is why the whole process is 
}


\section{Multithreaded Execution Backup and Restore}\label{sec:das-exec-restore}
When the program is executed, each section gets copied contiguously inside the virtual memory address range allocated for the program by the operating system. 
tried several techniques to save/restore execution stacks:
- sufficient conditions to be able to restore a task (w or w/o touching exec stack)
- different strategies, ultimately bounded by different reasons 
- strat 1 : just replay the thread until the first waitUntilNextPeriod() (there's 4 possible states : not created yet, created but not once had the "go", waiting in a waitUntilNextPeriod(), created but deleted)
- strat 2 : restart and stub functions, can do it with ELF-HOOK https://www.codeproject.com/Articles/70302/Redirecting-functions-in-shared-ELF-libraries
Redirection of the OSApiWrapper calls, Extremely useful, because our library is a dynamic object, which means its calls to library functions are always resolved at runtime by looking in the .rel.plt (relocation, Procedure Linkage Table)
- strat 3 complicated version of longjmp: how to save execution (asm, disassembly of code, save sp, etc.% https://cs.brown.edu/courses/cs033/docs/guides/x64_cheatsheet.pdf)
Depends on the ABI

}